% =====================================================================
%  c00.tex —— 标题页、底本说明、译者说明、目录（Lead 独占）
% =====================================================================
\begin{titlepage}
\centering
\vspace*{3.2cm}

{\Huge\bfseries 拓扑 $K$ 理论\par}
\vspace{0.8cm}
{\Large Robert Hines\par}
\vspace{0.5cm}
{\large 2016 年 12 月 15 日\par}

\vspace{2.2cm}
{\LARGE 中 文 译 注 本\par}
\vspace{0.6cm}
{\normalsize 底本：Robert Hines, \emph{Topological K-theory}, 2016-12-15（公开讲义）\par}

\vfill
{\small
\begin{tabular}{rl}
底本页数 & 10 页 \\
原文编号 & 定义 1--8；命题 1--3；引理 1；定理 1--9 \\
翻译体例 & 逐段对译 + 译者批注（不进入排版） \\
\end{tabular}\par}
\vspace{1.5cm}
{\small 本译本仅供个人学习研究使用；原文版权归原作者所有。\par}
\end{titlepage}

\pagenumbering{roman}

% ---------- 译者说明 ----------
\thispagestyle{empty}
\section*{译者说明}
\addcontentsline{toc}{section}{译者说明}

本译本是 Robert Hines 十页讲义 \emph{Topological K-theory}（2016-12-15）的中文译注本。原文以「$K(X)$ 由 $X$ 上的向量丛决定」为主线，依次交代向量丛与 $K(X)$、分类空间、
Bott 周期性，最后由 $\Omega$-预谱把 $K$ 理论放进广义上同调的框架，并以
「只有 $S^0,S^1,S^3,S^7$ 是 $H$-空间」为例说明 Adams 运算的用处。

体例说明：
\begin{itemize}[leftmargin=1.6em,itemsep=2pt]
  \item 正文为忠实逐段翻译，不增不减；原文笔误照译，仅在批注中指出并给出正确形式。
  \item 原文没有编号公式（实测 0 条），故译本亦无；原文的编号定义/命题/引理/定理一律手工给号，与原文一一对应。
  \item 原文的展示式（含矩阵、分段函数、交换图）全部按原文逐符号重排；交换图用 Ti\emph{k}Z 重画。
  \item 译者批注写成源码中的 \verb|% [Tutor Note: …]| 注释，不进入排版，另见随源码发布的《译注汇总.md》。
  \item 术语首现处括注英文原词；人名与书名保留原文。
\end{itemize}

\vspace{0.6em}
\noindent 页码对应：原文 PDF 页号 $=$ 印刷页码（偏移 0），正文中每起来自新一页处均有
\verb|% [原文 p.N]| 注释可对照。

\clearpage
\tableofcontents
\clearpage
